% Capturas de tela da aplicação cliente, organizadas por contexto de acesso.
% Cada tela: descrição + requisitos evidenciados + figura + legenda.
% As imagens devem ser salvas em docs/imagens/ com os nomes indicados.
% Enquanto a captura não existir, um marcador é exibido no lugar (o documento compila).
% IMPORTANTE: usar dados fictícios; não exibir tokens, e-mails reais, chaves ou IDs sensíveis.

\providecommand{\pendentePrint}[1]{%
  \fbox{\begin{minipage}[c][5.5cm][c]{0.55\textwidth}\centering
    \textit{[Captura de tela pendente]}\\[0.4cm]\texttt{imagens/#1}
  \end{minipage}}}

\providecommand{\telaPrint}[3]{% #1 = arquivo, #2 = largura, #3 não usado
  \IfFileExists{imagens/#1}{\includegraphics[width=#2\textwidth]{imagens/#1}}{\pendentePrint{#1}}}

As capturas a seguir ilustram as principais telas da aplicação cliente, todas com \textbf{dados fictícios de demonstração}. Para cada tela, indicam-se os requisitos funcionais por ela evidenciados.

\subsubsection{Catálogo público de viagens}

Tela inicial do contexto público, acessível sem autenticação. Apresenta as viagens disponíveis agrupadas por rota, com busca e filtros (data, turno e ordenação), e exibe a ocupação de cada viagem. Evidencia os requisitos \textbf{RF04} (viagens agendadas), \textbf{RF07} (ponto de entrada para inscrição) e \textbf{RF22} (indicação de capacidade/ocupação).

\begin{figure}[H]
\centering
\telaPrint{fe-catalogo.png}{0.55}{}
\caption{Catálogo público de viagens disponíveis.}
\label{fig:fe-catalogo}
\end{figure}

\subsubsection{Detalhe público da viagem}

Tela única de detalhe de uma viagem, funcional logada e deslogada. Mostra rota, horário, pontos de embarque e datas alternativas da mesma rota, além do botão de compartilhamento e do botão de inscrição (ou de \textit{login}, quando não autenticado). Evidencia \textbf{RF04} e \textbf{RF07}.

\begin{figure}[H]
\centering
\telaPrint{fe-detalhe-viagem.png}{0.55}{}
\caption{Detalhe público de uma viagem, com datas alternativas e compartilhamento.}
\label{fig:fe-detalhe-viagem}
\end{figure}

\subsubsection{Formulário de inscrição}

Formulário em que o passageiro autenticado escolhe os pontos de embarque e desembarque, o tipo de trajeto e a forma de pagamento. A decisão definitiva de disponibilidade e preço é delegada à API. Evidencia \textbf{RF07} (inscrição), \textbf{RF09} (tipo de trajeto), \textbf{RF10} (forma de pagamento) e \textbf{RF24} (seleção do ponto de embarque).

\begin{figure}[H]
\centering
\telaPrint{fe-formulario.png}{0.55}{}
\caption{Formulário de inscrição em uma viagem.}
\label{fig:fe-formulario}
\end{figure}

\subsubsection{Confirmação da inscrição}

Tela de confirmação exibida após a criação bem-sucedida da inscrição, evidenciando a conclusão do fluxo de reserva (\textbf{RF07}) e oferecendo acesso ao detalhe da inscrição.

\begin{figure}[H]
\centering
\telaPrint{fe-sucesso.png}{0.55}{}
\caption{Confirmação de inscrição concluída.}
\label{fig:fe-sucesso}
\end{figure}

\subsubsection{Minhas inscrições}

Lista das inscrições do passageiro, com busca e filtro por estado, distinguindo viagens futuras das já realizadas (histórico). Evidencia \textbf{RF08} (cancelamento dentro do prazo) e o acompanhamento do histórico de inscrições.

\begin{figure}[H]
\centering
\telaPrint{fe-minhas-inscricoes.png}{0.55}{}
\caption{Lista de inscrições do passageiro.}
\label{fig:fe-minhas-inscricoes}
\end{figure}

\subsubsection{Detalhe da inscrição}

Detalhe de uma inscrição específica, com os dados da viagem e a ação de cancelamento (sujeita à janela de prazo validada pela API). Evidencia \textbf{RF08}.

\begin{figure}[H]
\centering
\telaPrint{fe-detalhe-inscricao.png}{0.55}{}
\caption{Detalhe de uma inscrição, com opção de cancelamento.}
\label{fig:fe-detalhe-inscricao}
\end{figure}

\subsubsection{Painel administrativo}

Painel com métricas operacionais (viagens ativas, próximas, passageiros, ocupação) e financeiras (receita prevista), além de atalho para o relatório financeiro. Evidencia parcialmente \textbf{RF15} e \textbf{RF16} (relatório/histórico agregados no cliente) e \textbf{RF19} (estimativa financeira de arrecadação).

\begin{figure}[H]
\centering
\telaPrint{fe-dashboard.png}{0.75}{}
\caption{Painel administrativo com métricas operacionais e financeiras.}
\label{fig:fe-dashboard}
\end{figure}

\subsubsection{Gestão de viagens}

Tela administrativa de \textit{CRUD} de instâncias de viagem, com transições de estado e atribuição de motorista e veículo. Evidencia \textbf{RF04} (agendamento), \textbf{RF05} (cancelamento de viagem), \textbf{RF06} (turnos), \textbf{RF21} (vinculação de motorista/veículo) e \textbf{RF23} (pontos de embarque).

\begin{figure}[H]
\centering
\telaPrint{fe-gestao-viagens.png}{0.75}{}
\caption{Gestão administrativa de viagens.}
\label{fig:fe-gestao-viagens}
\end{figure}

\subsubsection{Gestão de templates de rota}

Tela de \textit{CRUD} de modelos de rota recorrentes e de geração manual de instâncias. Evidencia \textbf{RF04} (base do agendamento recorrente), \textbf{RF06} (turnos) e \textbf{RF23} (pontos de embarque do modelo).

\begin{figure}[H]
\centering
\telaPrint{fe-gestao-templates.png}{0.75}{}
\caption{Gestão de templates de rota.}
\label{fig:fe-gestao-templates}
\end{figure}

\subsubsection{Detalhe operacional da viagem}

Tela compartilhada por administrador e motorista para a gestão operacional das inscrições: confirmação de presença e de pagamento por passageiro e transições de estado da viagem. Evidencia \textbf{RF11} (lista de passageiros), \textbf{RF12} (confirmação de presença), \textbf{RF13} (lista para o motorista), \textbf{RF14} (início/conclusão da viagem) e \textbf{RF25} (pontos de embarque por passageiro).

\begin{figure}[H]
\centering
\telaPrint{fe-operacional.png}{0.75}{}
\caption{Detalhe operacional da viagem, com gestão de presença e pagamento.}
\label{fig:fe-operacional}
\end{figure}

\subsubsection{Gestão de motoristas e veículos}

Telas administrativas de gestão da frota: motoristas vinculados à organização e veículos cadastrados. Evidenciam \textbf{RF03} (cadastro de motoristas) e \textbf{RF20} (cadastro de meios de transporte).

\begin{figure}[H]
\centering
\telaPrint{fe-motoristas.png}{0.45}{}
\hfill
\telaPrint{fe-veiculos.png}{0.45}{}
\caption{Gestão de motoristas (à esquerda) e de veículos (à direita).}
\label{fig:fe-frota}
\end{figure}

\subsubsection{Viagens do motorista}

Lista das viagens atribuídas ao motorista autenticado, ponto de entrada para a tela de detalhe operacional. Evidencia \textbf{RF13} e \textbf{RF14}.

\begin{figure}[H]
\centering
\telaPrint{fe-motorista-viagens.png}{0.55}{}
\caption{Lista de viagens atribuídas ao motorista.}
\label{fig:fe-motorista-viagens}
\end{figure}
